\section{Scrum selection}

For this project, our goal was to create a heat optimization desktop application for Danfoss that could read data sheets, calculate production costs, track emissions and optimize based on two scenarios. Although these requirements were to a great extent well-understood from the outset of the project, we had to navigate technical uncertainty and our limited experience delivering such requirements in teams. Hence, Scrum was our preferred choice, even though it was strictly mandated by the semester project course and expectations, it enabled us to deliver iteratively, reflect on our experience and limitations, and adapt as we progressed through.

\section{Sprints}

We split our project timeline into a setup phase(Sprint 0), followed by five development sprints. If a sprint is too short, like one week, you spend too much time in meetings planning and reviewing instead of actually writing code. If a sprint is too long, like a month, we procrastinate because the deadline feels too far away. We kept each sprint two weeks long and found that a two-week schedule was the perfect length for our group as it gave us enough time to actually design, code, and test a working part of the app, but it was short enough to keep everyone focused and accountable.

\section{Team organization and roles}

In the spirit of Agile principles, we made a conscious decision to rotate the Scrum role among team members and emphasize collaboration over strict hierarchical structures and predefined roles. For example, throughout the project, the role of the Scrum Master was assumed by two different members. Instead of locking one person into an unchangeable position for the entire semester, this rotation strategy distributed ownership across the group. The Scrum Master focused on maintaining our Jira board, tracking deadlines and organizing our examinations of the code. This approach ensured that the responsibility for managing team workflows was shared fairlyand kept our sprints moving forward efficiently.

For the Product Owner (PO) role, we decided to share its responsibilities, which enabled a shared sense of accountability for the product's features and quality. Because our student group consists of six colleagues who all interacted directly with stakeholders and supervisors, isolating a single individual as our source of requirements would have slowed down communication. Instead the feedback we received from Danfoss in the midterm evaluation was owned by all team members. For example, when requested to add a CO2 emissions tracking graph, the entire group collaborated to map out the required features. Every member took part in translating these expectations into clear backlog items.

We used a volunteering approach in our meetings for asigning of the taks, where everybody volunteered based on their interests, technical strengths, and to keep everything fair, which worked really well for us and helped divide our work evenly.

\section{Scrum ceremonies adoption}

We followed the main ideas for the four Scrum meetings (Planning, Daily Standups, Reviews, and Retrospectives) but adapted them to fit our schedules as students. Even though we all shared the same university schedule, of course everyone also had commitments outside of classes and other personal activities.

Our Sprint Planning meetings took place every Monday at the start of the sprint and usually lasted around an hour. In this sessions we reviewed our backlog (see Appendix \ref{app:product_backlog}) and agreed on what we wanted to get done next. We used story points as a simple way of estimating the effort needed to finish them, but we also relied on our own inutition of how difficult a task seemed and how much free time we had available during that sprint. Before a task could be added to the sprint board everyone needed to understand what was required and how it should be implemented.

The Daily Scrum was the ceremony that we adapted the most. In theory Scrum suggests a short daily meeting where everyone shares updates. But for us, meeting in person every day was not the best use of our time. What we did instead was that we held one dedicated in person meeting each week to discuss progress and if any issues appeared and needed attention but for the rest of the week we stayed in contact through a WhatsApp group chat. Everyone posted updates whenever they completed a task or maybe if they ran into a problem which allowed us to stay informed without needing to meet every day and offeres us more flexibility.

We also combined the Sprint Review and Sprint Retrospective into a single session that took place when the sprint ended. In the review we tested the latest version of the app together and checked if the features worked as expected. After that we discussed how the sprint had gone, what worked well and what could be improved in the next sprint.

\section{Jira for workflow management}

When we set up the project in Sprint 0 we considered both Trello and Jira~\cite{jira}. Even though Trello was quite easy to understand and had a simple layout we thought Jira seemed better suited for a software development project because it offered more ways to organize our tasks and track the progress. One more reason is that we preferred Jira's interface because it is a very feature rich tool and the benefit of using it was that it is very used in the industry so getting some experience with it felt reasonable for future projects, therefore it was an easy choice.

Our setup was very simple, we organized the board into six columns that replicated how the tasks actually moved through the project:

\begin{itemize}
    \item Backlog - all of the existing tasks we have written (a comprehensive record of these items is available in Appendix \ref{app:product_backlog})
    \item To Do - tasks we chose for this sprint but haven't been started yet
    \item In Progress - tasks that were worked on in that moment
    \item In Review - tasks that were completed and for which the team needed to review the code first
    \item Testing - tasks that were completed but needed some manual testing to see if they behave correctly
    \item Done - tasks that have been approved by everybody and merged into the main branch~\cite{git}
\end{itemize}

Tasks often moved quickly between the columns because we were in constant connection with the whatsapp group, but having the board available made it easy to see what still needed attention before the end of each sprint.


\section{Delineations from Scrum recommendations}

It is clear that our team did not follow Scrum exactly as it is described but we treated it as a guideline and made it fit the realities of a six person student project.

Such framework is designed for professional development teams that work full-time but our situation was quite different and for us trying to follow every Scrum rule would probably just have added unnecessary complications without much additional value.

Because of this we focused on the parts of Scrum that genuinely helped us work together, as we still planned our work in sprints, maintained a backlog, held regular meetings and reflected on our process, but we adapted the details to suit our needs. For example we used story points as a rough guide instead of relying on time estimates, communicated through WhatsApp instead of holding daily in person standups and combined some meetings to spend less time spent on managing our collaboration. We deemed this lightweight implementation sufficient for our project and well-aligned with our academic circumstances.

Overall we believe our approach was successful even if it was less formal than a classic Scrum implementation, because it was well suited to our circumstances and most importantly it helped us deliver a fully functional application that met all of the requirements provided by Danfoss.
